\section{Parallelisierung} \label{sec:parallelisierung}
Um nun die Gipfelsuche auf mehreren Dateien zu beschleunigen kann die Gipfelsuche auf den Einzeldateien parallelisiert werden. \todo{vorschalg aus alter pa in text einarbeiten} Das würde bedeuten, dass alle Dateien gleichzeitig untersucht werden können und die Ergebnisse der einzelnen Dateien später wieder zusammengeführt werden können. Diese Architekturveränderung kann theoretisch die Geschwindigkeit bis zu der Anzahl der Prozessorkerne erhöhen. Im Weiteren werden verschiedene Ansätze der Paralellisierung besprochen und verglichen um den für diesen Anwendungsfall besten Ansatz zu finden.

Multithreading teilt ein Programm in mehrere Threads innerhalb eines Prozesses auf. Alle Threads teilen sich den gleichen Adressraum. Dadurch ist der Datenaustausch zwischen Threads effizient, allerdings steigt das Risiko von Race Conditions bei gemeinsam genutzten Datenstrukturen. Daher begrenzt der Global Interpreter Lock (GIL) in Python die gleichzeitige Ausführung von Python-Bytecode in mehreren Threads. Für gemeinsam genutzte Datenstrukturen sind zusätzlich explizite Sperren, etwa \textit{Lock}-Objekte, erforderlich, um konsistente Schreib- und Lesezugriffe sicherzustellen. 

Für die Verarbeitung mehrerer GeoTIFF-Dateien eignet sich Multithreading vor allem für I/O-lastige Teilschritte \todo{link zu ui reaktivität}, beispielsweise das Einlesen von Dateien. Bei CPU-lastigen Berechnungen ist der Geschwindigkeitsgewinn in Python (speziell in der Standard-Implementierung CPython) durch den GIL stark limitiert. Der GIL stellt sicher, dass stets nur ein einziger Thread zur selben Zeit den Python-Bytecode ausführt. Folglich kann ein einzelner Python-Prozess selbst mit mehreren aktiven Threads immer nur maximal einen CPU-Kern physisch auslasten. Darum handelt es sich bei CPU-lastigen Aufgaben in Python um kein „echtes“, parallel ausgeführtes Multithreading über mehrere Kerne, sondern lediglich um ein schnelles, nebenläufiges Wechseln (Nebenläufigkeit) auf einem Kern. Dieser Aspekt sorgt dafür, dass der classische Multithreadingansatz in Python nicht zu einer Perfromancesteigerung führt. Daher muss bei dem Gipfelsuchalgorithmus auf andere Verfahren zurückgegriffen werden.

Als Alternative kann Multiprocessing eingesetzt werden. Dabei wird jede Datei in einem separat gestarteten Prozess analysiert, von denen jeder seinen eigenen Interpreter samt eigenem GIL besitzt. Nur dieser Ansatz erlaubt in Python eine echte Parallelität auf mehreren Prozessorkernen, erfordert jedoch mehr Speicher und einen höheren Verwaltungsaufwand für die Prozesskommunikation. Dieser Ansatz hat einen entscheidenen Nachteil, denn er sorgt für eine Hohe Ramauslastung. Dies hängt damit zusammen, dass Python in jedem Prozess alle Daten halten muss die in diesem Prozess verwendet werden. Diese Tatsache sorgt bei der parallelen Ausführung mehrerer Dateien (nach \ref{sec:multi_file_execution}) für den doppelten Ramverbrauch im Gegensatz zur sequenzieler Ausführung. Bei der Zusammenführung (nach \ref{sec:multi_file_merge}) sorgt das für eine unmögliche Ausführung des Algorihtmus. Das liegt daran, dass hier für jeden Gipfel, der untersucht werden muss alle GeoTIFF-Dateien gespeichert werden müssen. Das bedeutet der Ramverbrauch ist bei der Zusammenführung um ein vielfaches höher als zuvor. Dies führt in den meisten fällen zu einem Absturz des Computers.

Um Multiprocessing nun rameffizienter umzusetzen gibt es drei Möglichkeiten. Die erste wäre eine Gruppierung der Gipfel. Dies meint, dass die Gipfel oder Dateien vor der Untersuchtung zuerst in Gruppen aufgeteilt werden und danach für jede der Gruppen ein Prozess erstellt wird, der die Analyse für diese ausführt. Wenn die Anzahl der Gruppen immer dem \textit{max(Prozessorkernen, Analyseobjekten)} entspricht wird durch diese Änderung die Geschwindigkeit des Multiprocessing nahezu beibehalten. Er ist lediglich wegen dem Austausch der Analysepbjekte nahc jedem Durchlauf langsamer. Der Ramverbrauch ist außerdem geringer. Er ist je nach Gerät und GeoTIFF-Dateigröße immernoch im kritischen Bereich und kann zum Einfrieren oder Abstürzen des Computer führen. 

Eine weitere Möglichkeit wäre es anstatt bei der Erstellung eines prozesses alle Daten zu übergeben, jedem Prozess immer nur die Daten zu geben, wenn er diese benötigt. Dies würde dazu führen, dass dauerhaft weniger Speicher beansprucht wird, würde aber die Ausführzeit erhöhen. Zudem wäre dies sehr aufwändig zu implementieren. Alternativ hier zu können die Daten alle anstatt an den Prozess selbst übergeben zu werden in einem gemeinsamen Speicher (shared Memory) liegen, auf den der jeder Prozess selbst zugreift. Dies würde die Ausführgeschwindigkeit des Algorithmus zwar minimal verringern, aber den Ramverbrauch enorm senken.

Für den vorliegenden Anwendungsfall ist daher eine differenzierte Strategie zweckmäßig: I/O-lastige Schritte werden mittels Multithreading parallelisiert, CPU-lastige Analysen pro Datei mittels Multiprocessing mit shared Memory. Die Ergebnisse der Einzelanalysen werden anschließend zentral zusammengeführt und in \ref{sec:multi_file_merge} dateiübergreifend validiert. Diese validierung wird auch mit der Hilfe von Multiprozessing mit shared Memory parallelisiert.